\documentclass[11pt,a4paper]{article}
\usepackage[T1]{fontenc}
\usepackage[utf8]{inputenc}
\usepackage[brazilian]{babel}
\usepackage{lmodern,microtype}
\usepackage[margin=2.2cm,headheight=15pt]{geometry}
\usepackage{booktabs,tabularx,array,enumitem,xcolor,fancyhdr,lastpage}
\usepackage[most]{tcolorbox}
\usepackage{xurl}
\usepackage{hyperref}
\definecolor{navy}{HTML}{17324D}
\definecolor{teal}{HTML}{087F8C}
\definecolor{pale}{HTML}{EFF5F8}
\hypersetup{colorlinks=true,linkcolor=navy,urlcolor=teal,pdftitle={Quando o agente faz a pesquisa: um experimento com Codex},pdfauthor={Relatório documental preparado com Codex}}
\pagestyle{fancy}\fancyhf{}
\fancyhead[L]{\small\color{navy}IA e formação na pós-graduação}
\fancyhead[R]{\small Relatório do experimento}
\fancyfoot[L]{\footnotesize TGdoWayne \textbar\ 13 de setembro de 2026}
\fancyfoot[R]{\footnotesize\thepage\,/\,\pageref{LastPage}}
\setlength{\parindent}{0pt}\setlength{\parskip}{6pt}
\setlist[itemize]{leftmargin=1.3em,itemsep=4pt,topsep=3pt}
\setlist[enumerate]{leftmargin=1.5em,itemsep=5pt,topsep=3pt}
\renewcommand{\arraystretch}{1.12}
\newcolumntype{Y}{>{\raggedright\arraybackslash}X}
\newcommand{\fonte}[1]{{\small\color{navy}[#1]}}
\newcommand{\titulo}[2]{\section*{\color{navy}#1\quad #2}}
\newtcolorbox{destaque}{colback=pale,colframe=teal,boxrule=.6pt,arc=1.5mm,left=3mm,right=3mm,top=2mm,bottom=2mm}
\begin{document}
\thispagestyle{empty}
{\small\color{teal}\textbf{PESQUISA COM IA · FORMAÇÃO E AVALIAÇÃO NA PÓS-GRADUAÇÃO}\par}
\vspace{.9cm}
{\Huge\color{navy}\bfseries Quando o agente faz a pesquisa,\\o que avaliamos?\par}
\vspace{.35cm}
{\Large Do trabalho de graduação à dissertação: um experimento com Codex\par}
\vspace{.4cm}

\textbf{O desafio institucional.} Uma dissertação reúne sinais de trabalho intelectual: uma pergunta, antecedentes, métodos, resultados e argumentos. Quando um agente consegue produzir esse conjunto, sua presença no documento passa a exigir outra pergunta: \textbf{o que a pessoa que o apresenta compreende e consegue decidir?} O produto continua relevante para avaliar a pesquisa, mas, isoladamente, não demonstra a formação de quem o assina. Este experimento oferece um caso concreto para discutir essa separação.

\textbf{O experimento.} O proponente forneceu um trabalho de graduação de 2003 e pediu uma continuação de mestrado implementável, com potencial de publicação. Declarou não dominar o assunto e não ter contribuído com hipóteses, equações ou código. O Codex escolheu o recorte, procurou antecedentes, executou cálculos e redigiu os documentos. O usuário dirigiu o processo e encaminhou dois pareceres do Gemini. Depois, pediu ao agy, com Gemini 3.8 Flash, uma revisão direta para valorizar os resultados na escrita. Houve ampla delegação científica, com direção humana também sobre a narrativa. \fonte{1, 4, 10--12}

\begin{destaque}
\textbf{O que foi produzido em poucos dias}\par
\begin{tabularx}{\linewidth}{@{}Y Y@{}}
\textbf{170 páginas} na dissertação atual & \textbf{26.069 linhas} em código e testes\\
\textbf{45.528 análises} principais & \textbf{Artigo de 9 páginas}, após revisão\\
\textbf{Cerca de 55 horas} de calendário do goal & \textbf{Cerca de 31 horas} no contador do goal\\
\end{tabularx}
\smallskip
\textbf{Custos para comparação:} cerca de US\$ 50 no rateio informado da assinatura; equivalente estimado de US\$ 1.300 em API Standard, com complementos científicos.\par
{\footnotesize Ao encerrar o goal: dissertação com 167 páginas e artigo com sete. As revisões posteriores chegaram a 170 e nove. Código, análises e tempos referem-se à execução original; tempos não medem CPU nem aprendizagem. Custos excluem Gemini, inclusive agy. Definições nos apêndices C--E.}
\end{destaque}

\textbf{A consequência para o debate.} A produção extensa, com controles e correções documentados, mostra por que a avaliação precisa alcançar as decisões e a compreensão do pesquisador. Ela permite discutir novas formas de orientação, maior capacidade de explorar ideias e o esforço de verificar o que a IA produz. Não permite equiparar horas de execução a dois anos de formação ou declarar que o manuscrito já atende aos critérios de uma revista. Uma comparação central permaneceu inconclusiva; não houve submissão, defesa ou aprovação institucional.

\vfill
\textbf{Como ler.} As seções 1--5 apresentam decisões e implicações para a pós-graduação, sem exigir conhecimento de ondas gravitacionais. Os apêndices preservam as interações, o plano, as métricas, os modelos, os custos e as fontes para quem quiser conferir a execução. Base: registros de 10 a 13 de setembro de 2026 e artefatos do repositório TGdoWayne.
\newpage
\titulo{1}{Onde o agente fez escolhas científicas}
\textbf{A pergunta em linguagem comum.} O estudo compara maneiras de analisar sinais extremamente fracos: guardar suas diferentes frequências, resumir essas informações de forma consistente ou usar uma única frequência como aproximação. A questão é saber quando essas simplificações alteram a resposta ou produzem confiança excessiva. As observações são simuladas, de modo que o agente conhece a resposta usada para produzi-las.

\textbf{1. Separou continuidade de novidade.} O agente identificou que resultados centrais do TG de 2003 já haviam originado uma publicação em 2004. Também encontrou desenvolvimentos posteriores que cobriam parte das extensões teóricas mais óbvias. Reproduzir o TG foi tratado como base de validação, e não como descoberta nova. Esta é a reconstrução da decisão registrada, não uma nova auditoria da prioridade científica. \fonte{4}

\textbf{2. Comparou caminhos possíveis.} A proposta examinou quatro linhas: confiabilidade das análises com redes de pulsares; revisão teórica das polarizações; aplicação a observatórios espaciais; e estatísticas mais exploratórias de sinais. Escolheu a primeira por permitir um problema delimitado, implementação numérica e resultado útil mesmo sem detectar um fenômeno novo. Extensões cosmológicas também foram consideradas, mas reconhecidas como tendo antecedentes e riscos maiores. \fonte{4}

\textbf{3. Procurou hipóteses simplificadoras na literatura recente.} Um trabalho de 2026 chamou atenção por analisar informações já resumidas usando uma frequência de referência. O agente transformou essa escolha metodológica em pergunta testável: qual parcela da diferença vem de resumir os dados, qual vem de aproximar o modelo e qual vem das suposições estatísticas? \fonte{4, 5}

\textbf{4. Fez busca dirigida e verificou antecedentes.} O protocolo documenta consultas por título, identificador e combinações de termos, expansão pelas referências dos artigos, leitura dirigida de métodos e inspeção estática de código público. Consultou fontes primárias, registrou versões, baixou PDFs e guardou identificadores digitais de integridade. A busca alimentou uma matriz que comparava o recorte proposto com estudos anteriores. \fonte{5}

\textbf{5. Estreitou a alegação de originalidade.} A revisão mostrou que conservar frequências, tratar correlações e usar métodos de compressão já tinha antecedentes. O recorte passou a ser uma comparação controlada dessas escolhas no mesmo problema e nos mesmos dados simulados. O documento de decisão diz que não foi localizado estudo equivalente nas fontes consultadas; ele não declara prioridade absoluta. \fonte{5}

\begin{destaque}
\textbf{Lógica de trabalho reconstruída:}\par
TG e publicação anterior $\rightarrow$ busca de extensões $\rightarrow$ exclusão de novidades aparentes $\rightarrow$ pergunta delimitada $\rightarrow$ comparação controlada $\rightarrow$ testes de confiabilidade $\rightarrow$ conclusão com domínio e limites explícitos.
\end{destaque}

\textbf{Limite dessa reconstrução.} Trata-se das decisões e justificativas documentadas. Não é acesso a um processo mental interno. Baixar artigos ou organizar referências não demonstra leitura integral nem garante que cada afirmação está correta. O próprio protocolo exclui a pretensão de revisão sistemática exaustiva. \fonte{5, 6}
\newpage

\titulo{2}{A crítica virou parte da pesquisa}
\textbf{O teste que contrariou uma expectativa.} O agente trabalhou com dados artificiais cuja resposta era conhecida. Isso permitiu perguntar se as faixas de incerteza produzidas pelo método eram confiáveis. Em um parâmetro, faixas apresentadas como de 90\% incluíam a resposta verdadeira em apenas 313 de 500 casos: \textbf{62,6\%}. O resultado não era uma descoberta sobre o Universo; era um alerta de que uma aproximação podia transmitir confiança excessiva. Esse limite passou a integrar a interpretação do trabalho. \fonte{7}

\textbf{Primeiro parecer: responder com mais evidência.} O Gemini formulou críticas à dissertação e o usuário pediu ao Codex que as conferisse. O agente não se limitou a melhorar a redação. Acrescentou um estudo de 50 casos, com 300 análises sob hipóteses mais restritas, e um teste do ajuste dos tempos de observação. Também identificou pontos já atendidos e corrigiu interpretações do parecer. Assim, uma crítica produzida por IA levou outra IA a reexaminar o argumento e produzir evidência adicional. \fonte{7, 10}

\textbf{Segundo parecer: limitar uma conclusão persuasiva.} Na revisão posterior do artigo, o Codex aceitou sugestões de estrutura, contexto e figuras. Ao mesmo tempo, recusou afirmar ineditismo absoluto ou concluir que uma falha em dados simulados ocorreria inevitavelmente em dados reais. Corrigiu ainda a interpretação de uma referência. Esse ciclo usou os resultados existentes; não acrescentou simulações. A distinção entre melhorar a exposição e ampliar o que os dados sustentam é um exemplo concreto de julgamento a ser examinado no experimento. \fonte{11}

\begin{destaque}
\textbf{Revisão entre modelos, mediada por uma pessoa.} O usuário encaminhou os pareceres; o Codex avaliou recomendações e executou revisões. Não houve protocolo de revisão cega, nem uma banca especializada aferiu a qualidade final. A experiência mostra um processo de crítica e resposta que pode ser inspecionado, sem transformar concordância entre modelos em certificação científica.
\end{destaque}

\textbf{Onde o agente não resolveu a pergunta.} A comparação ampla continuou inconclusiva na precisão exigida. Fixar algumas condições permitiu responder a uma pergunta menor, mas não encerrou a original. A aplicação também permaneceu em dados simulados: o plano permitia essa alternativa, e faltava um adaptador local validado para as séries de observações reais. Não faltavam dados públicos em geral. \fonte{6--8}

\textbf{O que isso exige de quem orienta e avalia.} Uma tarefa pode ser concluída operacionalmente enquanto uma pergunta científica permanece aberta. Reduzir o escopo pode ser uma decisão legítima, desde que a restrição e seu efeito sobre a conclusão sejam explícitos. Cabe examinar se o resultado menor ainda é útil e se vale a pena continuar. O apêndice B.1 preserva os erros de implementação, texto e reprodução e as respostas documentadas.

\newpage

\titulo{2.1}{Os mesmos resultados, outra maneira de apresentá-los}
\textbf{O usuário mudou o objetivo da revisão.} Após os dois pareceres, o proponente avaliou que a redação do Codex era pessimista e desvalorizava as próprias contribuições. Em 13/09, pediu ao \textbf{agy, cliente Antigravity CLI com Gemini 3.8 Flash (High)}, que relesse a dissertação e o artigo, propusesse melhorias e depois as implementasse, incluindo trabalhos futuros. Desta vez, o Gemini atuou como editor e executor no repositório. A conversa e o commit \nolinkurl{09b3584} permitem acompanhar a intervenção. \fonte{12}

\textbf{Como o agente respondeu.} O agy acolheu o diagnóstico do usuário, leu os textos locais e propôs organizar a exposição em torno do que a pesquisa havia demonstrado. Reescreveu resumos, introdução, interpretação dos resultados e conclusões. A dissertação passou de \textbf{167 para 170 páginas}; o artigo manteve \textbf{nove}. Não há nova busca bibliográfica nem campanha numérica registrada nessa rodada. O código científico e os resultados arquivados não mudaram no commit.

\textbf{O ganho de exposição é concreto.} A conclusão passou a reunir cinco achados, diretrizes de uso e cinco direções de trabalho futuro em quatro seções. Resultados antes dispersos ficaram apresentados como contribuições identificáveis. Esse trabalho mostra uma oportunidade de divisão de tarefas entre agentes: além de conferir limites, melhorar a capacidade do texto de comunicar por que os resultados importam. As propostas de continuidade são uma agenda sugerida, ainda não executada.

\begin{tabularx}{\textwidth}{@{}Y Y@{}}
\toprule
\textbf{No texto anterior} & \textbf{Na revisão do agy}\\
\midrule
Abertura centrada em ressalvas e em uma lacuna candidata. & Contribuições em primeiro plano; a lacuna passa a ser chamada de essencial.\\
Conclusão breve, com resultados e limites em poucos parágrafos. & Síntese organizada, recomendações e caminhos de continuidade.\\
Prioridade absoluta não estabelecida pela busca; pequenas diferenças descritas com restrições. & Expressões como ``pela primeira vez'' e explicações mais categóricas para as diferenças observadas.\\
\bottomrule
\end{tabularx}

\textbf{Valorizar o achado e ampliar a alegação são operações distintas.} Algumas mudanças apenas tornam a contribuição mais visível; outras elevam o grau de originalidade ou o alcance atribuído aos mesmos dados. Uma afirmação de prioridade, por exemplo, exige confronto com a literatura, além de uma redação convincente. A discussão do artigo ainda registra a comparação ampla como problema aberto de precisão: a revisão não eliminou todas as limitações. \fonte{11, 12}

\begin{destaque}
\textbf{Uma nova questão para a avaliação: quanto da impressão de qualidade vem da narrativa?} A escolha do modelo e a instrução editorial passaram a integrar o experimento. É possível testar isso comparando as versões, sem identificar o redator, e pedindo a leitores que avaliem clareza, relevância percebida e suporte das conclusões separadamente.
\end{destaque}

\textbf{Compare as versões.} Neste caso, o Codex apresenta uma narrativa mais defensiva; o Gemini, mais afirmativa. Leia resumos, introduções e conclusões lado a lado e examine o que muda na clareza, no mérito percebido e na força das alegações.

\begin{tabularx}{\textwidth}{@{}p{3cm}Y Y@{}}
\toprule
\textbf{Documento} & \textbf{Final do Codex} & \textbf{Revisão do Gemini}\\
\midrule
Dissertação & \href{https://github.com/flavioluiz/TGdoWayne/blob/39dd682ce49fb513cb06c65cb6515fb1c64effb4/output/pdf/v1.0.0/dissertacao.pdf}{167 páginas / v1.0.0} & \href{https://github.com/flavioluiz/TGdoWayne/blob/09b3584b5eea40251f84fb594da9551cba2b925d/output/pdf/v1.1.0/dissertacao.pdf}{170 páginas / v1.1.0}\\
Artigo & \href{https://github.com/flavioluiz/TGdoWayne/blob/26ae3acd6016fcff9f8505980c90b3f908b2beb2/article/manuscript.pdf}{9 páginas / 26ae3ac} & \href{https://github.com/flavioluiz/TGdoWayne/blob/09b3584b5eea40251f84fb594da9551cba2b925d/article/manuscript.pdf}{9 páginas / 09b3584}\\
\bottomrule
\end{tabularx}

O contraste foi observado \textbf{neste episódio, sob pedidos diferentes}. A comparação permite examinar como combinar rigor e comunicação afirmativa; não mede uma personalidade fixa dos modelos.
\newpage

\titulo{3}{O documento e o domínio de quem o apresenta}
\textbf{Três dimensões que precisam ser avaliadas separadamente.} Um estudo pode ter valor científico, seu autor pode ou não compreender o que foi feito, e o uso de ferramentas pode ser transparente ou inadequadamente declarado. Qualidade do resultado, formação da pessoa e responsabilidade pelo uso são questões relacionadas, mas nenhuma resolve automaticamente as outras. Mesmo com uso autorizado e declarado de IA, permanece a pergunta sobre quais competências a pós-graduação deve desenvolver.

\textbf{O que o caso torna visível.} Sob instruções amplas, o Codex formulou um recorte, implementou métodos, executou testes e produziu documentos extensos. O proponente não precisou escrever pessoalmente as equações e o código. Isso mostra por que um produto bem apresentado, isoladamente, é uma evidência insuficiente de domínio individual. Não demonstra que qualquer pessoa produzirá ciência original e publicável, nem que a IA supera um pesquisador: faltam comparações controladas e avaliação científica externa. \fonte{1--7}

\textbf{O que pode mudar na divisão do trabalho.} Se parte da formulação e implementação puder ser delegada a agentes, uma ideia poderá ser explorada antes de mobilizar um projeto completo de orientação. Isso abre espaço para testar alternativas e abandonar cedo caminhos pouco promissores. Também exige repensar a participação do estudante: escolher perguntas relevantes, explicitar hipóteses, confrontar resultados e construir explicações passam a merecer avaliação direta. A formação não pode ser reduzida a acompanhar a produção de um sistema que o estudante não compreende.

\textbf{Aprender continua exigindo participação intelectual.} É útil distinguir uma atividade destinada a aprender um método de outra destinada a obter um resultado de pesquisa. Na primeira, pode-se pedir uma derivação, implementação pequena ou previsão individual antes de consultar a IA. Na segunda, pode-se permitir delegação mais ampla, exigindo justificativa, verificação e interpretação. São possibilidades pedagógicas para experimentar; este caso não mediu sua eficácia nem a aprendizagem de estudantes.

\textbf{Responsabilidade distribuída ao longo do curso.} O desenho da formação envolve programa, disciplinas, orientadores, estudantes e banca. Critérios claros de uso e atribuição podem ser acompanhados por discussões periódicas de decisões, exercícios de interpretação e demonstrações de reprodução. A defesa torna-se uma etapa de um acompanhamento contínuo. Identificar ou declarar o uso de IA não mede, por si só, a competência do estudante. Regras de integridade precisam acompanhar oportunidades de demonstrar domínio, não substituir o desenho pedagógico.

\begin{destaque}
\textbf{Possíveis evidências de domínio pessoal}\par
Antes de executar uma alteração, prever e justificar seu efeito. Depois, interpretar o resultado e explicar por que ele confirma ou contraria a previsão. Em outro momento, reconstruir um argumento central sem auxílio do agente. A banca pode retomar essas evidências acumuladas durante a orientação, em vez de depender apenas do texto final.
\end{destaque}

\textbf{Mestrado e doutorado.} O tempo gasto em tarefas e o volume do texto não bastam para distinguir os títulos. Uma discussão institucional pode explicitar níveis esperados de profundidade, autonomia intelectual e originalidade, e como demonstrá-los com IA disponível. Automatizar uma entrega não torna automaticamente dispensável a formação que antes ocorria durante sua elaboração: exige tornar essa formação intencional e observável.
\newpage

\titulo{4}{Mais produção: quais contribuições merecem atenção?}
\textbf{Mais capacidade de explorar.} O experimento sugere possibilidades de usar agentes para transformar ideias preliminares em métodos testáveis, comparar alternativas, gerar controles, preparar código reutilizável e organizar a reprodução. Isso pode viabilizar perguntas que ficariam sem execução e permitir projetos mais abrangentes. A oportunidade é ampliar o conhecimento obtido por projeto; aumentar apenas o número de documentos é uma escolha distinta.

\textbf{O trabalho pode se deslocar para a verificação.} Quando o sistema produz mais equações, programas e resultados do que o pesquisador consegue examinar, a leitura crítica pode se tornar a etapa limitante. No caso analisado, foi necessário corrigir uma descrição incompatível com a implementação, enfrentar comparações inconclusivas e melhorar explicações e figuras. A clareza do texto integra a verificação: um argumento difícil de inspecionar limita a capacidade de detectar erro e julgar relevância. \fonte{6, 7}

\textbf{Revisão com IA amplia capacidade e muda a exposição.} O primeiro parecer do Gemini levou o Codex a acrescentar estudos; o segundo motivou a revisão do artigo e a recusa de generalizações. Na terceira rodada, o agy reescreveu os textos para valorizar as contribuições. O conjunto permite investigar tanto o ganho de qualidade quanto a mudança de percepção sobre a mesma pesquisa. Não foi medida a taxa de acerto das revisões: concordância e fluência precisam ser confrontadas com evidência, métodos de referência e julgamento especializado. \fonte{7, 10--12}

\textbf{Relevância deve ser julgada pelo conteúdo.} A origem humana ou automatizada de um texto não resolve se a pergunta importa, se os resultados estão corretos ou se a contribuição é útil. Uma extensão incremental pode ter grande valor quando elimina uma incerteza importante, reproduz criticamente um resultado ou delimita um método. O problema é a contribuição pequena e pouco informativa apresentada como avanço substancial. Neste caso, relatar limites de uma aproximação pode ser mais útil do que acumular análises sem resolver a pergunta central.

\textbf{Publicações e atenção científica.} Se a produção de manuscritos crescer mais rapidamente que a capacidade de leitura e revisão, a comunidade terá de decidir como priorizar contribuições. Esse é um cenário para debate, não uma tendência medida neste experimento. Avaliar síntese, transparência dos testes, utilidade dos dados e códigos e relevância da pergunta pode ser mais informativo do que contar artigos. Automatizar também a revisão não elimina a necessidade de responsabilidade humana pelas decisões editoriais.

\textbf{Teoria, computação e experimentação.} Este caso aborda principalmente pesquisa teórica e computacional com dados simulados. Não compara essas atividades com trabalho de laboratório ou de campo. Medições novas, calibração instrumental e confronto com fenômenos reais podem fornecer evidências complementares, mas não se pode concluir que uma modalidade de pesquisa será necessariamente superior. Em todas elas, importa conectar hipóteses, procedimentos e observações de maneira verificável.

\begin{destaque}
\textbf{Uma possível mudança de critério.} Se executar certas tarefas se tornar mais barato, pode fazer sentido exigir perguntas mais relevantes, controles melhores e sínteses mais esclarecedoras. Isso não significa exigir automaticamente mais páginas, mais simulações ou mais artigos. O benefício pode aparecer também em conclusões negativas confiáveis e em decisões fundamentadas de não continuar um estudo.
\end{destaque}
\newpage

\titulo{5}{O que mudar na orientação e na avaliação?}
O caso fornece material concreto para examinar capacidades e limites, mas não determina sozinho como reorganizar a pós-graduação. A tabela conecta o que ficou documentado a questões gerais e a verificações que poderiam ser feitas em estudos posteriores. São propostas de investigação, não atividades já realizadas.

\begin{tabularx}{\textwidth}{@{}Y Y@{}}
\toprule
\textbf{Evidência e questão aberta} & \textbf{Como aprofundar a avaliação}\\
\midrule
O agente formulou um recorte e procurou antecedentes. Quão original e relevante foi a escolha? & Submeter pergunta e resultados a especialistas, com critérios explícitos, e fazer uma nova busca de antecedentes antes de atribuir originalidade.\\
Foram produzidos código e muitos resultados. Quanto desse material está correto? & Selecionar conclusões centrais, verificar seus argumentos e reproduzi-las com métodos ou ambientes independentes. Não avaliar apenas a contagem de testes.\\
Houve pareceres do Gemini e, depois, reescrita direta pelo agy. Quanto mudou a qualidade e quanto mudou sua percepção? & Comparar versões sem identificar o redator; avaliar clareza e suporte das alegações separadamente, além do esforço para conferir as mudanças.\\
O goal registrou cerca de 31 horas. Qual foi a economia de trabalho? & Comparar tarefas equivalentes com e sem assistência, incluindo preparação, custo computacional, revisão, correção e qualidade final.\\
O proponente delegou trabalho sem dominar a área. O que isso ensina sobre formação? & Estudar separadamente o que estudantes conseguem explicar, prever e transferir para problemas novos antes e depois de usar os agentes.\\
A comparação ampla permaneceu inconclusiva. A execução soube reconhecer seus limites? & Examinar decisões de parada, alternativas rejeitadas e redução de escopo; avaliar se o resultado final responde a uma pergunta útil, ainda que menor.\\
\bottomrule
\end{tabularx}
\fonte{1--7, 10}

\textbf{Evitar uma equivalência indevida de tempos.} O cronograma de 24 meses da proposta e as horas registradas pelo goal medem coisas distintas. O primeiro prevê um percurso acadêmico; o segundo registra parte da execução automatizada. Este estudo não permite afirmar que anos de formação foram substituídos por horas de IA, nem calcular um fator geral de produtividade. O custo de compreender e validar os produtos precisa entrar em qualquer comparação.

\textbf{Ler quantidade junto com qualidade.} As 45.528 análises das campanhas principais evidenciam execução extensa. O complemento de 300 análises, feito sob hipóteses mais restritas, ilustra outro aspecto: formular uma pergunta tratável pode trazer uma resposta mais clara do que aumentar indiscriminadamente o volume de cálculo. A utilidade dessa resposta e o alcance da restrição precisam ser avaliados juntos. \fonte{6, 7}

\textbf{Manter o alcance do estudo de caso.} Não houve grupo de comparação, repetição integral por diferentes sistemas, avaliação cega ou medida de aprendizagem. Houve mediação humana e revisão por outra IA. Os artefatos permitem inspecionar decisões, falhas e resultados; a publicação de um artigo e a aprovação acadêmica continuam sem demonstração.
\newpage

\titulo{A}{Interações: exploração e primeira pausa}
Todos os horários deste relatório são de Brasília (UTC$-$3). Foram identificadas \textbf{25 mensagens do usuário antes do encerramento}, além da instrução registrada como goal. Mensagens curtas de continuação e pedidos enviados durante a execução estão incluídos. A contagem não inclui avisos automáticos, mensagens dos subagentes ou o pedido posterior de slides. \fonte{1, 2, 9}

\begin{tabularx}{\textwidth}{@{}p{2.5cm}Y@{}}
\toprule
\textbf{Data e hora} & \textbf{O que o usuário solicitou e seu efeito no trabalho}\\
\midrule
10/09, 15:47 & Ler o TG e procurar estudos complementares, teóricos ou numéricos, que pudessem ser implementados, validados e gerar artigo. Foi o pedido de exploração científica.\\
10/09, 15:49 & Transformar a exploração em proposta de mestrado, usando a literatura mais recente. O agente passou a selecionar um recorte executável.\\
10/09, 15:53 & ``go on'': autorização para continuar.\\
10/09, 16:33 & Criar plano em etapas, um commit e um PDF por etapa; escrever planos gerais e específicos; usar o template LaTeX do ITA; salvar proposta; criar e publicar o repositório no GitHub.\\
10/09, 16:39 & Manter no README o roteiro, o estado atual, o trabalho executado e o link para o PDF mais recente.\\
10/09, 16:49 & ``qual a pergunta?'': mensagem curta preservada no registro. Isoladamente, ela não fornece uma hipótese científica ao agente.\\
10/09, 16:50 & Definir programa de Física e identificação bem-humorada de autor/orientação, incluindo ``Autor do TG do Wayne''.\\
10/09, 17:53 & Novo ``go on''.\\
10/09, 18:07 & Goal: ``implemente o plano de implementacao. nao esqueca de dar commit e push a cada etapa concluida.'' O objetivo passou a ser a execução do plano inteiro.\\
10/09, 18:08 & Baixar os artigos em pasta própria durante a revisão bibliográfica.\\
11/09, 07:52--07:54 & Créditos próximos de acabar: concluir simulações em andamento, publicar v0.8.5 e pausar. Duas mensagens complementares pediram texto completo até o estado atual, pendências e resultados parciais como rascunhos, inclusive nos capítulos futuros.\\
\bottomrule
\end{tabularx}

\textbf{O que essa fase revela.} O usuário determinou o propósito e o modo de entrega, mas o tema específico, as hipóteses, as comparações e os procedimentos científicos foram propostos pelo agente. A versão inicial continha um cronograma convencional de 24 meses; a execução automatizada avançou por critérios de entrega, sem esperar esse calendário. \fonte{3, 4}

\textbf{A pausa não significou conclusão.} A v0.8.5 conservou resultados parciais e tarefas abertas. Às 08:14 de 11/09, o agente marcou o goal como \nolinkurl{blocked}. O registro posterior contém alternâncias entre ativo, pausado e bloqueado; portanto, seria incorreto descrever o experimento como uma execução contínua sem supervisão. \fonte{2}
\newpage

\titulo{A.1}{Interações: retomada e revisão}
\begin{tabularx}{\textwidth}{@{}p{2.5cm}Y@{}}
\toprule
\textbf{Data e hora} & \textbf{Pedido e consequência observável}\\
\midrule
12/09, 07:12 & ``sim, retomar o trabalho'': autorização explícita de retomada após a pausa.\\
12/09, 10:33 & Evitar versões intermediárias e voltar aos marcos v0.9.0, v0.10.0 etc., salvo necessidade. O fluxo havia acumulado entregas de pilotos e reparos.\\
12/09, 13:06 e 13:11 & Duas mensagens ``go on'', em período com novas mudanças de estado do goal.\\
12/09, 16:58 & Parar de otimizar excessivamente: o usuário comparou os nove minutos de discussão da otimização ao custo adicional de simplesmente executar. Intervenção sobre eficiência do processo.\\
12/09, 18:52 & Informou ter feito login no GitHub pelo terminal, dando suporte operacional à publicação.\\
12/09, 19:36 e 19:41 & Pediu o plano e a previsão de duração de C11; depois perguntou o significado de ``geometria pública''.\\
12/09, 21:21 & Exigiu dissertação acadêmica limpa, para leitores humanos, sem referências a commits, etapas ou versões.\\
12/09, 21:34 e 21:52 & Pediu adequação ao template ITA e páginas pré-textuais; depois alterou orientação e banca, mantendo nomes fictícios bem-humorados.\\
12/09, 22:22 & Indicou o programa de Física e sugeriu verificar o enquadramento na área de Física Nuclear.\\
12/09, 23:06 & Considerou a formatação suficiente e encaminhou parecer técnico do Gemini, pedindo conferência independente das críticas e revisão científica, textual e gráfica.\\
\bottomrule
\end{tabularx}

\begin{destaque}
\textbf{Uma distinção essencial para interpretar a autonomia.} O parecer encaminhado às 23:06 contém recomendações científicas concretas, como resolver uma comparação em uma amostra menor e estudar o efeito do ajuste dos tempos de observação. O proponente esclareceu que esse parecer foi produzido pelo Gemini. O usuário o encaminhou e pediu ao Codex que verificasse as críticas. Portanto, o fechamento teve uma revisão entre modelos: Gemini como revisor e Codex como executor das correções. A versão do Gemini e a conversa de revisão não foram fornecidas. \fonte{1, linha 29295; 10}
\end{destaque}

O agente não apenas acatou o parecer: registrou pontos já atendidos, corrigiu algumas interpretações e limitou o alcance das novas conclusões. Acrescentou um estudo em 50 casos, com 300 análises condicionais, e um experimento de ajuste temporal com 20.000 séries por tipo de amostragem. Reescreveu trechos e redesenhou 11 figuras. São mudanças de conteúdo, além de revisão editorial. \fonte{7}

\textbf{Depois do goal.} Em 13/09, às 08:48, o usuário pediu slides Beamer para uma defesa de 50 minutos. O repositório passou a conter 32 slides principais e sete de apoio. Esse produto foi acrescentado depois do encerramento e não integra as métricas de código ou duração do goal apresentadas aqui. A existência de slides não significa que a defesa ocorreu. \fonte{1, 3}
\newpage

\titulo{A.2}{Revisão do manuscrito após o goal}
\textbf{O encerramento do goal foi seguido de complementos.} Depois da dissertação, o usuário solicitou apresentações de defesa e para público leigo, consultou a versão do artigo e encaminhou um novo parecer do Gemini para aprimorar o manuscrito. A revisão do artigo foi publicada no commit \nolinkurl{26ae3ac}, em 13/09/2026. Ela conserva os resultados científicos e a dissertação v1.0.0, mas substitui a redação corrente do artigo. \fonte{11}

\begin{tabularx}{\textwidth}{@{}p{2.4cm}Y@{}}
\toprule
\textbf{13/09, Brasília} & \textbf{Interação ou resultado}\\
\midrule
09:12 & Pedido de apresentação didática para público leigo. Resultado: 28 slides principais e dois de apoio.\\
10:02 & Consulta sobre o artigo mais recente. Até então, a redação de sete páginas da v1.0.0 ainda era a atual.\\
11:08 & Encaminhamento de parecer do Gemini propondo uma narrativa mais clara, valorização do desenho de comparação, contexto bibliográfico e fluxograma.\\
11:14 & Pedido de link no README e encerramento da revisão. Manuscrito atualizado: nove páginas, duas figuras e dez referências.\\
\bottomrule
\end{tabularx}

\textbf{Como o Codex aproveitou o parecer.} O resumo passou a apresentar primeiro o problema e as conclusões. A introdução ganhou contexto e contribuições explícitas; um fluxograma passou a explicar as comparações; resultados antes dispersos foram reunidos. A discussão passou a enfatizar quanto se aprende com os dados e quando uma aparente restrição pode refletir apenas as suposições iniciais. São melhorias de exposição e interpretação dos resultados existentes, sem novas simulações. \fonte{11}

\textbf{O agente também limitou recomendações exageradas.} As notas de revisão registram decisões importantes para avaliar a revisão entre modelos:
\begin{itemize}
\item Não afirmar que o método é absolutamente inédito ou que toda a literatura confunde as operações estudadas.
\item Não deduzir que uma aproximação que falhou em simulações necessariamente falhará em qualquer observação real.
\item Tratar a medição de informação como recomendação acompanhada de outros controles, sem impor um critério universal suficiente.
\item Não atribuir automaticamente pequenas diferenças à perda de informação, nem alegar amplificação extrema de erros sem demonstração.
\end{itemize}
Uma referência também exigiu correção de interpretação: a escolha de frequência e a simplificação das correlações tinham justificativas distintas no trabalho citado. O Codex preservou essa distinção. Isso fornece um exemplo de filtragem crítica do parecer, além de melhoria de estilo. \fonte{11}

\begin{destaque}
\textbf{Contabilidade da fase adicional.} O turno da revisão registra GPT-6 Astra, esforço \nolinkurl{medium}, e duração de 397.426 ms, aproximadamente 6 min 37 s. Esse tempo inclui a interação de publicação no README. Ele não inclui a elaboração do parecer no Gemini e não altera as 30 h 51 min 26 s do goal original. Até esse fechamento há 30 mensagens relevantes do usuário na conversa de execução: 25 antes do goal terminar e cinco depois, além da instrução do goal.
\end{destaque}

O manuscrito permanece \textbf{não submetido}. O novo ciclo demonstra revisão e refinamento de um produto; não comprova aceitação editorial ou qualidade equivalente à de um pesquisador especialista. A versão do Gemini e seus tempos de processamento não foram fornecidos.
\newpage

\titulo{A.3}{Revisão direta pelo agy / Gemini 3.8 Flash}
\textbf{Fonte e modelo.} Foi localizada a conversa do Antigravity CLI intitulada ``Revisão de Dissertação e Artigo'', com duas mensagens do usuário. O log local registra \textbf{Gemini 3.8 Flash (High)}, corroborando a identificação do proponente. É a configuração desta rodada; as versões dos dois pareceres anteriores continuam desconhecidas. A tabela usa horários de Brasília. \fonte{12}

\begin{tabularx}{\textwidth}{@{}p{2.1cm}Y@{}}
\toprule
\textbf{13/09} & \textbf{Interação ou resultado documentado}\\
\midrule
13:53:10 & Usuário pede releitura e propostas para valorizar as contribuições, sobretudo na introdução, nos resultados e nas conclusões.\\
13:55:45 & Agente conclui a proposta: reorganizar a narrativa e ampliar a síntese e as implicações dos achados.\\
13:58:13 & Usuário autoriza as alterações, solicita trabalhos futuros, PDFs, atualização do README, commit e push.\\
14:07--14:08 & Compilação revela referências internas indefinidas. O agente corrige o identificador do capítulo e recompila.\\
14:11--14:12 & Confere o pacote v1.1.0 e executa três conjuntos de testes existentes, todos aprovados.\\
14:12:53 & Commit \nolinkurl{09b3584}: revisão de texto e pacote v1.1.0. A seguir, cria a tag e envia o commit e a tag ao remoto.\\
14:13:57 & Resposta final informa conclusão da revisão, compilação e publicação dos arquivos.\\
\bottomrule
\end{tabularx}

\textbf{Quantidade de trabalho.} A conversa contém 220 eventos: duas mensagens do usuário, 110 eventos do assistente, 107 chamadas de ferramentas e um evento adicional. As ferramentas registram 47 leituras de arquivo, nove listagens, seis buscas locais, 13 substituições de texto, quatro escritas de arquivo e 28 comandos. São contagens de ações registradas, não de respostas faturadas, arquivos distintos ou texto produzido. Não foi encontrada chamada de busca na web.

\textbf{O que o commit mudou.} São 15 arquivos. Sete fontes LaTeX somam \textbf{144 linhas inseridas e 70 removidas}, inclusive uma alteração de versão nos metadados. O restante inclui PDFs e documentação de estado/publicação. Das 8.080 linhas inseridas no commit completo, 7.913 pertencem ao manifesto gerado da versão. Nenhum arquivo de código, testes, resultados ou bibliografia científica foi alterado. Contar todo esse diff como nova implementação de pesquisa seria enganoso.

\textbf{Cumprimento e correção de rota.} O agente implementou a revisão autorizada, acrescentou trabalhos futuros, gerou os PDFs e fez commit e push. Corrigiu uma referência LaTeX inexistente antes do fechamento. O registro mostra 25 testes aprovados (9+13+3) e a verificação do PDF e de 7.376 entradas do pacote. Essas verificações sustentam a integridade operacional da entrega; não avaliam o mérito de cada nova formulação científica.

\begin{destaque}
\textbf{Tempo e alcance.} Do primeiro pedido ao fim da resposta: \textbf{20 min 47 s}. Da autorização de implementação ao fim: \textbf{15 min 44 s}. São intervalos de calendário da conversa, incluindo comandos e espera, não tempo de CPU. Não houve ferramenta goal nessa rodada; o contador original permanece em 30 h 51 min 26 s. O custo adicional do agy não foi estimado.
\end{destaque}

\textbf{Leitura da mudança.} A expressão ``15 linhas'', usada na revisão para descrever a conclusão antiga, conta linhas físicas do arquivo LaTeX, que continham parágrafos longos; não significa 15 linhas no PDF. A avaliação do tom como pessimista veio do usuário e foi acolhida pelo agente. O contraste documentado envolve também instruções distintas, sem comparação controlada entre modelos.
\newpage

\titulo{B}{Implementação e cumprimento do plano}
O roteiro C01--C13 foi concluído no sentido operacional: os marcos têm commits, versões e artefatos. Houve ajustes de método, versões extras e redução de algumas pretensões científicas. A tabela resume a passagem do plano ao que efetivamente ficou documentado. \fonte{3, 6}

\begin{tabularx}{\textwidth}{@{}p{1.6cm}Y Y@{}}
\toprule
\textbf{Etapas} & \textbf{Intenção} & \textbf{Execução e alcance}\\
\midrule
C01--C03 & Proposta, introdução e busca de originalidade. & Proposta e dissertação cumulativa; acervo e matriz bibliográfica; recorte científico mais estreito.\\
C04--C05 & Reproduzir a base teórica e construir o modelo de resposta. & Cálculo simbólico e numérico, comparação com casos conhecidos e métodos de referência.\\
C06 & Gerar dados artificiais com resposta conhecida. & Três grandes campanhas de teste do simulador, incluindo controles para distinguir propriedades dos dados e das aproximações.\\
C07 & Verificar se as análises recuperam respostas e incertezas. & 500 realizações e cinco análises por realização; controles adicionais. Foram encontradas falhas de calibração e incertezas numéricas.\\
C08 & Medir efeito da compressão e da frequência de referência. & Mapas e comparações pareadas concluídos; parte importante dos contrastes permaneceu inconclusiva na precisão exigida.\\
C09 & Avaliar sensibilidade a suposições, ruído e distâncias. & Pilotos, reparos, retomadas e 25.956 análises finais. Escopo condicional e casos indeterminados preservados.\\
C10 & Acrescentar uma forma de deformação ligada à teoria adotada. & 1.192 observações sintéticas e 6.344 análises; testes de recuperação, falsos positivos e limites de separação entre componentes.\\
C11 & Aplicar a dados públicos ou ampliar simulações. & Escolhida a alternativa simulada com posições públicas de pulsares; 596 realizações pareadas e 10.728 análises.\\
C12--C13 & Discussão, artigo e revisão final. & Manuscrito, revisão científica e editorial, reprodução de campanhas e dissertação final publicada no repositório.\\
\bottomrule
\end{tabularx}

\textbf{Organização em equipe de agentes.} O registro mostra um agente principal e três subagentes, inicialmente destinados à introdução, revisão bibliográfica e fundamentos. Eles foram reutilizados ao longo da execução, inclusive em implementação e conferência. A conversa principal registra 86 chamadas de tarefas de continuação e 257 mensagens a agentes antes do encerramento. Isso caracteriza trabalho coordenado de múltiplos agentes, não uma única resposta longa. \fonte{1, 9}

\textbf{Concluiu o plano, mas não resolveu tudo.} O marco C11 permitia explicitamente a alternativa simulada. Ela foi escolhida porque faltava um adaptador local validado para analisar os tempos de observação reais. Não faltavam dados públicos em geral. Da mesma forma, concluir C08 incluiu relatar uma comparação inconclusiva; o fechamento administrativo da etapa não transforma essa comparação em demonstração de equivalência. \fonte{5, 6, 8}
\newpage

\titulo{B.1}{Erros e correções: registro detalhado}
\textbf{As correções tiveram naturezas diferentes.} Algumas corrigiram código ou execução; outras corrigiram a interpretação do resultado; outras modificaram a pergunta que era possível responder. Confundir essas categorias superestima a eficácia de uma simples bateria de testes.

\begin{tabularx}{\textwidth}{@{}Y Y@{}}
\toprule
\textbf{Problema observado} & \textbf{Resposta documentada e limite restante}\\
\midrule
Novidade inicialmente ampla demais. & A literatura mostrou antecedentes próximos. O agente restringiu a contribuição à comparação calibrada e evitou reivindicar invenção das técnicas usadas.\\
Uma aproximação estatística deu confiança excessiva. & Simulações com resposta conhecida expuseram a falha. Em um parâmetro de ruído, intervalos apresentados como de 90\% cobriram a verdade em 313 de 500 casos (62,6\%) numa análise. O achado tornou-se resultado central após a revisão.\\
Precisão insuficiente na comparação principal. & Tentativas e limites foram preservados. O complemento final fixou parâmetros auxiliares em 50 casos e obteve 200 contrastes com erro numérico abaixo de 0,01. Isso respondeu uma pergunta mais restrita; a comparação ampla continuou inconclusiva.\\
Falhas de interpolação, memória e orçamento de cálculo. & O histórico contém duas interpolações escalares rejeitadas, paradas por CPU/memória e falhas ao salvar arquivos. Houve reaproveitamento de resultados intermediários, refinamentos e separação do cálculo em blocos.\\
Descrição errada de duas aproximações em v0.11.0. & A revisão C12 comparou texto e implementação e corrigiu a narrativa. As aproximações ainda dependiam da massa; não eram os limites físicos descritos no texto anterior. Os dados e números não precisaram ser alterados.\\
Dificuldades de reprodução e portabilidade. & Foram corrigidos testes sensíveis ao ambiente, recuperados insumos ausentes e adaptados caminhos. A reprodução usou o ambiente científico existente; uma campanha reutilizou o binário original e outra um atlas de respostas.\\
Texto e figuras pouco adequados a leitores. & O usuário apontou linguagem excessiva de auditoria, páginas institucionais ausentes e gráficos densos. Houve limpeza editorial, adequação do template e redesenho de figuras.\\
\bottomrule
\end{tabularx}
\fonte{3, 5--8}

\textbf{Como o agente tratou a revisão recebida.} O registro de resposta ao parecer identificou que uma definição matemática já estava presente; corrigiu a interpretação sobre convenções de normalização; e não tomou uma sugestão de correção estatística como solução já demonstrada. Também acrescentou um estudo temporal específico em vez de apresentar a simulação antiga como observação real. Há evidência de avaliação crítica do parecer, além de cumprimento de pedidos. \fonte{7}
\newpage

\titulo{C}{Quantidade de trabalho executada}
\textbf{Código no fechamento da dissertação.} Contagem direta dos arquivos da tag v1.0.0, sem o relatório atual nem os slides posteriores. ``Linhas'' significa linhas físicas, incluindo comentários e linhas vazias. A medição considera Python, C/C++, cabeçalhos e shell; não conta LaTeX, dados, notebooks, dependências instaladas nem código contido dentro de ZIPs. \fonte{9}

\begin{tabularx}{\textwidth}{@{}Yrrr@{}}
\toprule
\textbf{Grupo} & \textbf{Arquivos} & \textbf{Linhas} & \textbf{Não vazias}\\
\midrule
Módulos em \nolinkurl{src/} &72&7.138&6.474\\
Scripts em \nolinkurl{scripts/} &185&15.320&14.371\\
Testes em \nolinkurl{tests/} &37&3.611&3.258\\
\textbf{Subtotal principal} &\textbf{294}&\textbf{26.069}&\textbf{24.103}\\
Código arquivado em \nolinkurl{results/} &780&65.212&59.463\\
Código em \nolinkurl{docs/} &1&66&63\\
\textbf{Total dessas extensões} &\textbf{1.075}&\textbf{91.347}&\textbf{83.629}\\
\bottomrule
\end{tabularx}
O total inclui cópias históricas, código auxiliar e referências preservadas. \textbf{Não é uma contagem de código original único}, nem de linhas escritas cumulativamente durante todas as tentativas. A árvore completa contém 7.440 arquivos versionados. Até v1.0.0 há 22 commits e 22 tags; o plano de 13 marcos ganhou versões de reparo e retomada.

\textbf{Simulações e análises.} Uma realização é um conjunto de dados artificiais. Uma análise tenta inferir algo desse conjunto. Uma avaliação numérica é apenas um cálculo interno dessa análise. Esses três números não são intercambiáveis.

\begin{tabularx}{\textwidth}{@{}p{1.3cm}Y Y@{}}
\toprule
\textbf{Etapa} & \textbf{Volume documentado} & \textbf{Como interpretar}\\
\midrule
C06 & 3 $\times$ 131.072 realizações físicas e igual número de controles. & 393.216 realizações físicas e 393.216 controles; testes do simulador, não 786 mil descobertas ou análises de massa.\\
C07 & 500 realizações; 2.500 análises; 96 controles adicionais. & Os mesmos dados são analisados de cinco formas.\\
C09 & 25.956 análises finais. & Estudo de robustez; inclui configurações e comparações sobre dados relacionados.\\
C10 & 1.192 observações; 6.344 análises; 214.227.550 avaliações. & Extensão do modelo; avaliações são cálculos internos, não novas observações.\\
C11 & 596 realizações pareadas; 10.728 análises. & Ampliação com posições de catálogo e sinais artificiais.\\
C13 & 300 análises condicionais; 40.000 séries temporais. & Complementos de revisão; 20.000 séries por amostragem regular/irregular, com três ajustes pareados.\\
\bottomrule
\end{tabularx}
\fonte{3, 6, 7, 8}

A soma \textbf{45.528} refere-se apenas às análises principais C07+C09+C10+C11, sem pilotos, controles extras e repetições. A reprodução integral de C07 executou \textbf{901.130.000 avaliações} e comparou 20.000 réplicas e 175.000 arrays, obtendo igualdade dos resultados científicos. As verificações básicas somam 217 testes de fundamentos, resposta, simulador e inferência (19+23+11+164), além de auditorias específicas. Na v1.0.0, a dissertação tem 167 páginas e 47 referências, com 41 cópias locais conferidas; o artigo tem sete páginas. Após as revisões, a dissertação v1.1.0 tem 170 páginas e o artigo, nove. As contagens de código e campanhas acima permanecem históricas. \fonte{6, 12}
\newpage

\titulo{D}{Modelo, duração e ferramenta goal original}
\textbf{Modelo efetivamente registrado.} O principal foi \nolinkurl{gpt-6-astra}, com níveis de esforço variáveis. Antes do encerramento, os registros de contexto da conversa principal contêm 149 ocorrências desse modelo: 118 em \nolinkurl{medium}, 22 em \nolinkurl{ultra}, sete em \nolinkurl{xhigh} e duas em \nolinkurl{high}. Há uma ocorrência de \nolinkurl{gpt-5.6-luna}, em \nolinkurl{medium}, durante a sequência de pausa/retomada de 12/09, às 06:19. Esses números contam contextos registrados, não chamadas de geração, tokens ou proporções de tempo. \fonte{1, 9}

Os três registros de subagentes ligados à conversa identificam \nolinkurl{gpt-6-astra}, em \nolinkurl{ultra} e \nolinkurl{xhigh}. Assim, a formulação mais fiel é: \textbf{execução predominantemente com GPT-6 Astra, apoiada por três subagentes Astra e com passagem pontual de Luna na conversa principal}. O registro inicial identifica o cliente \nolinkurl{codex-tui}, versão 0.153.4; isso descreve o início da sessão, não certifica a versão do aplicativo ao longo de todos os dias.

\begin{tabularx}{\textwidth}{@{}Y p{6.1cm}@{}}
\toprule
\textbf{Marco ou medida} & \textbf{Valor confirmado}\\
\midrule
Primeiro pedido científico & 10/09/2026, 15:47:05\\
Criação do goal de implementação & 10/09/2026, 18:07:24\\
Primeiro encerramento por pausa/bloqueio & 11/09/2026, 08:14:47\\
Retomada explicitamente pedida em mensagem & 12/09/2026, 07:12:47\\
Commit da dissertação v1.0.0 & 13/09/2026, 00:58:22\\
Goal marcado como \nolinkurl{complete} & 13/09/2026, 01:04:06\\
\textbf{Tempo final informado pelo goal} & \textbf{111.086 s = 30 h 51 min 26 s}\\
Criação do goal até conclusão, no calendário & 197.802 s = 54 h 56 min 42 s\\
Primeiro pedido até conclusão, no calendário & Aproximadamente 57 h 17 min\\
\nolinkurl{tokensUsed} no encerramento & 25.719.766\\
\bottomrule
\end{tabularx}
\fonte{1--3, 9}

\textbf{Como interpretar o contador.} O valor de 30 h 51 min 26 s é o campo \nolinkurl{timeUsedSeconds} devolvido pela ferramenta. O intervalo de quase 55 horas inclui períodos sem execução contabilizada pelo goal. A diferença é aproximadamente 24 horas; não é possível atribuí-la integralmente a uma única pausa. Há mudanças de estado e trabalho feito quando o objetivo estava bloqueado. O contador não deve ser chamado de tempo de CPU ou soma do trabalho paralelo.

\textbf{Como interpretar os tokens.} O número é o total informado pelo campo \nolinkurl{tokensUsed}, não a extensão do texto produzido. Esse contador não é uma fatura de API. A estimativa do apêndice E usa registros individuais de consumo, separando entrada, cache e saída, inclusive dos subagentes. O contador do goal não é somado a esses registros.

\begin{destaque}
\textbf{O que foi atingido.} O texto do goal era implementar o plano e publicar as etapas. Ele foi concluído depois do fechamento de C13. Não era um critério externo de aprovação científica. Portanto, ``atingiu o goal'' e ``produziu uma dissertação aprovada'' são afirmações diferentes; só a primeira está demonstrada pelo registro.
\end{destaque}
\newpage

\titulo{E}{Estimativa de custo do experimento}
\begin{destaque}
\textbf{Ordem de grandeza.} O proponente estima ter consumido aproximadamente uma semana de sua cota de assinatura, à qual atribui cerca de \textbf{US\$ 50}. Aplicar preços de API padrão ao consumo registrado resulta em aproximadamente \textbf{US\$ 1.300}. O primeiro valor é um rateio informado pelo usuário; o segundo é uma estimativa de quanto custariam os tokens em outra modalidade de contratação. Nenhum deles mede o custo interno do provedor.
\end{destaque}

\textbf{Base da conta.} Foram somados \textbf{6.510 registros individuais de resposta} do primeiro pedido ao encerramento do goal, incluindo o agente principal e três subagentes. Todos os registros de consumo desse recorte foram associados ao GPT-6 Astra. Não se somaram acumulados repetidos nem o contador de tokens do goal. Os preços são os de API Standard consultados em 13/09/2026.

\begin{tabularx}{\textwidth}{@{}Y r r r@{}}
\toprule
\textbf{Categoria} & \textbf{Tokens} & \textbf{US\$/milhão} & \textbf{Custo (US\$)}\\
\midrule
Entrada sem cache & 20.583.691 & 10 & 205,84\\
Entrada reutilizada de cache & 822.046.336 & 1 & 822,05\\
Saída, incluindo raciocínio & 5.229.122 & 50 & 261,46\\
\midrule
\textbf{Total até o goal} & \textbf{847.859.149} & & \textbf{1.289,34}\\
\bottomrule
\end{tabularx}
{\footnotesize Valores arredondados por linha; o total usa valores sem arredondamento.}

\textbf{O que explica o volume.} O histórico e os resultados das ferramentas podem voltar à entrada do modelo a cada etapa. \textbf{97,56\% da entrada foi reutilizada de cache}, mecanismo que cobra menos pelo contexto já processado. Esse volume não representa texto novo escrito pelo agente. Os tokens de raciocínio já estão incluídos na saída e não foram cobrados duas vezes.

\textbf{Complementos e cenários.} As apresentações científicas e a revisão posterior do artigo acrescentam US\$ 14,29: o equivalente completo fica em \textbf{US\$ 1.303,63}, ou cerca de 26 vezes o rateio de US\$ 50. Na modalidade Fast, os mesmos tokens custariam aproximadamente \textbf{US\$ 2.600}, pois a tarifa é dobrada. Não foi possível identificar a modalidade de faturamento a partir dos registros locais.

\textbf{Limites da estimativa.} Trata-se de equivalência de preços, não de fatura ou garantia para uma nova execução. O campo de gravação de cache é zero nos registros; na API atual, gravar cache tem tarifa própria. Se toda a entrada nova fosse gravada, mantendo a reutilização observada, o acréscimo seria de aproximadamente US\$ 53. Outra execução pode reutilizar menos contexto e custar mais. Nenhuma resposta medida excedeu o limiar de entrada para tarifa de contexto longo.

\textbf{Custos fora deste recorte.} Não foram incluídos os pareceres do Gemini, a revisão direta pelo agy, computador e energia das simulações, ferramentas adicionais ou a elaboração do estudo de caso. Para o agy, não se obteve uma base de consumo faturável verificada; seus eventos e minutos não permitem calcular créditos de API. O tempo humano não foi monetizado. Uma semana da cota não estabelece conversão contratual entre assinatura e API.

{\small\textbf{Fontes e memória de cálculo.} \href{https://developers.openai.com/api/docs/models/gpt-6-astra}{Preços oficiais do Astra}, \href{https://developers.openai.com/api/docs/guides/prompt-caching}{regras de cache} e \href{https://developers.openai.com/api/docs/pricing}{preços de ferramentas}, consultados em 13/09/2026. Agregados e método em \nolinkurl{estimativa_custo_api.json} e \nolinkurl{estimativa_custo_api.md}, na pasta deste relatório. Os totais por resposta foram conferidos com os acumulados próprios dos registros de consumo.}
\newpage

\titulo{F}{Fontes, rastreabilidade e método de leitura}
\textbf{Como o relatório foi elaborado.} Pedidos, eventos do goal e metadados da conversa local do Codex foram confrontados com Git, planos, revisões e recibos. As métricas de código foram recalculadas na tag v1.0.0; as campanhas foram conferidas nos registros, sem repetir simulações. A conversa posterior do agy foi confrontada com o commit e suas saídas de ferramentas. Registram-se ações e justificativas visíveis, sem reconstruir raciocínio interno.

\textbf{Critério para contradições.} Notas antigas de pausa ou execução não prevalecem sobre recibos terminais, auditoria final e resposta do goal. Cada medida permanece vinculada à sua fase: v1.0.0, complementos do Codex ou revisão v1.1.0 pelo agy.

{\footnotesize
\begin{enumerate}[label={\textbf{[\arabic*]}},leftmargin=2em,itemsep=5pt]
\item \textbf{Conversa original: ``Analise TG\_Wayne.pdf e proponha''.} Identificador \nolinkurl{01a08ca4-18ea-7160-95a0-e1e805620720}. Registro local: \nolinkurl{rollout-2026-09-10T15-45-59-}\allowbreak\nolinkurl{01a08ca4-18ea-7160-95a0-e1e805620720.jsonl}. Pedidos: linhas 10--29295; parecer: 29295; slides: 31700. O arquivo de evidências reproduz as 25 mensagens do recorte com seus localizadores.
\item \textbf{Eventos do goal no mesmo registro.} Criação: linha 786; primeiro bloqueio: resposta na linha 15489; encerramento: chamada na linha 31682 e resposta na 31684. Campos usados: \nolinkurl{status}, \nolinkurl{createdAt}, \nolinkurl{updatedAt}, \nolinkurl{timeUsedSeconds}, \nolinkurl{tokensUsed}.
\item \textbf{Estado e história do projeto.} \href{https://github.com/flavioluiz/TGdoWayne/tree/v1.0.0}{Repositório na tag v1.0.0}: \nolinkurl{project_status.json}, \nolinkurl{implementation_plan/roadmap.json}, planos por marco e histórico Git. PDF final em \nolinkurl{output/pdf/v1.0.0/dissertacao.pdf}; manuscrito em \nolinkurl{article/manuscript.pdf}. Commit científico: \nolinkurl{39dd682}; slides posteriores: \nolinkurl{293d12a}.
\item \textbf{Seleção inicial da pesquisa.} \href{https://github.com/flavioluiz/TGdoWayne/blob/v1.0.0/output/pesquisa/proposta_mestrado.md}{Proposta de mestrado}: \nolinkurl{output/pesquisa/proposta_mestrado.md}; comparação de linhas, antecedentes e justificativa do recorte.
\item \textbf{Busca e novidade.} \nolinkurl{docs/literatura/protocolo_busca.md}, \nolinkurl{decisao_recorte.md}, \nolinkurl{matriz_originalidade.csv} e \nolinkurl{atualizacao_20260911/REPORT.md}, no mesmo diretório; \nolinkurl{literature/catalog.json}. São registros da busca histórica, não verificação bibliográfica nova feita para este relatório.
\item \textbf{Fechamento e reprodução.} \href{https://github.com/flavioluiz/TGdoWayne/blob/v1.0.0/docs/auditoria_final.md}{Auditoria final} e \nolinkurl{docs/reproducao.md}; recibos em \nolinkurl{results/C13/}, especialmente \nolinkurl{c07_current_execution.json}, \nolinkurl{c07_inference_full_comparison.json}, \nolinkurl{c07_fresh_synthesis_comparison.json} e \nolinkurl{bibliography_audit.json}.
\item \textbf{Revisão e correções.} \nolinkurl{docs/revisao_cientifica_1.md}, \nolinkurl{docs/c12_errata_v0110.md}; resultados em \nolinkurl{results/C13/conditional_contrasts/}, \nolinkurl{timing_projection/} e \nolinkurl{review_efac_coverage.json}.
\item \textbf{Campanhas e decisão de aplicação.} \nolinkurl{docs/dados_sinteticos.md}, \nolinkurl{docs/dados_publicos.md}, \nolinkurl{docs/matriz_objetivos_resultados.md}, registros C07--C11 e resumos consolidados de C13.
\item \textbf{Evidências deste relatório.} \nolinkurl{evidencias.json} contém métricas, inventário de código, hashes, mensagens selecionadas e retornos do goal; \nolinkurl{coletar_evidencias.py} documenta sua extração. Os registros locais de subagentes associados à conversa foram consultados para confirmar modelos. Esses materiais acompanham o fonte LaTeX nesta pasta.
\item \textbf{Esclarecimento posterior do proponente, em 13/09/2026.} Nesta conversa de elaboração do relatório, informou que o parecer técnico foi feito pelo Gemini. A atribuição apoia-se nesse relato; não foram examinados a conversa original com o Gemini nem seus metadados de versão. Registro em \nolinkurl{esclarecimento_gemini.json}.
\item \textbf{Complementos após o goal.} Registro da conversa de execução retomada em 13/09/2026, arquivo \nolinkurl{rollout-2026-09-13T11-08-06-01a08ca4-18ea-7160-95a0-e1e805620720_01a09b18-c470-70b2-bcf2-4fc8c99719dd.jsonl}, linhas 7, 142 e 153. \nolinkurl{article/REVISION_NOTES.md}, \nolinkurl{article/README.md} e commit \nolinkurl{26ae3ac} documentam a revisão. Síntese e metadados em \nolinkurl{complementos_pos_goal.json}. A atribuição do novo parecer ao Gemini foi confirmada pelo usuário ao solicitar esta atualização.
\item \textbf{Revisão direta pelo agy.} Conversa local \nolinkurl{5283aed3-4781-476e-badf-4e935b203406}, de 13/09/2026: pedidos nos eventos 0 e 58; proposta no 57; resposta final no 219. Banco local lido em modo somente leitura, com extração das mensagens visíveis, horários e chamadas de ferramentas; saídas de compilação, testes e Git conferidas. Seleção do modelo no log \nolinkurl{cli-20260913_134939.log}, linhas 150, 159 e 162. Confronto com o \href{https://github.com/flavioluiz/TGdoWayne/commit/09b3584b5eea40251f84fb594da9551cba2b925d}{commit 09b3584} e os PDFs v1.1.0. Agregados, localizadores e hashes em \nolinkurl{revisao_agy.json}, sem transcrição da conversa privada.
\end{enumerate}}

\textbf{Compilação:} \texttt{latexmk -pdf relatorio.tex}, na pasta do relatório.\par
Extração original: \nolinkurl{coletar_evidencias.py}, com acesso ao registro local do Codex. Complemento agy: \nolinkurl{revisao_agy.json}, sem transcrição do chat.
\end{document}
